% ============================================================
%  Relatório: Detecção de Fraudes em Transações Financeiras
%  Métodos Baseados em Árvores
% ============================================================
\documentclass[12pt,a4paper]{article}

% ── Pacotes ──────────────────────────────────────────────────────────────────
\usepackage[utf8]{inputenc}
\usepackage[T1]{fontenc}
\usepackage[brazil]{babel}
\usepackage{geometry}
\geometry{margin=2.5cm}

\usepackage{amsmath,amssymb}
\usepackage{booktabs}
\usepackage{multirow}
\usepackage{graphicx}
\usepackage{float}
\usepackage{caption}
\usepackage{subcaption}
\usepackage{hyperref}
\hypersetup{colorlinks=true, linkcolor=blue, citecolor=blue, urlcolor=blue}
\usepackage{xcolor}
\usepackage{microtype}
\usepackage{enumitem}
\usepackage{array}
\usepackage{siunitx}
\sisetup{group-separator={.}, group-minimum-digits=4}

% ── Macros ───────────────────────────────────────────────────────────────────
\newcommand{\dataset}[1]{\textsc{#1}}
\newcommand{\algo}[1]{\textit{#1}}
\newcommand{\metric}[1]{\texttt{#1}}

% ── Metadados ────────────────────────────────────────────────────────────────
\title{%
  \textbf{Detecção de Fraudes em Transações Financeiras}\\[6pt]
  \large Comparação de Métodos Baseados em Árvores de Decisão
}
\author{Lucas Pacheco}
\date{Março de 2026}

% ═════════════════════════════════════════════════════════════════════════════
\begin{document}

\maketitle
\tableofcontents
\newpage

% ─────────────────────────────────────────────────────────────────────────────
\begin{abstract}
Este relatório apresenta um estudo comparativo de seis métodos de detecção de
fraudes em transações financeiras, todos pertencentes à família de métodos
baseados em árvores de decisão: \algo{XGBoost}, \algo{LightGBM},
\algo{CatBoost}, \algo{Isolation Forest}, \algo{Stacking} e
\algo{Super Ensemble}. Os experimentos foram conduzidos em dois \textit{benchmarks}
amplamente utilizados na literatura — o dataset \dataset{Credit Card Fraud
Detection} (CC) e o \dataset{IEEE-CIS Fraud Detection} (IEEE) — que diferem
substancialmente em dimensionalidade, taxa de fraude e tipo de features
disponíveis. Os resultados evidenciam que métodos de \textit{gradient boosting}
supervisionados (XGBoost, CatBoost) dominam na maioria das métricas, enquanto
o \algo{Isolation Forest} — único método não supervisionado avaliado — apresenta
ROC-AUC competitivo mas PR-AUC muito inferior, reflexo direto do desbalanceamento
severo das classes. O \algo{Super Ensemble} obtém as melhores métricas no
dataset IEEE (PR-AUC = 0{,}5095; F1 = 0{,}5255), ao custo de tempo de
treinamento consideravelmente maior.
\end{abstract}

% =============================================================================
\section{Introdução}
\label{sec:intro}
% =============================================================================

A detecção de fraudes em pagamentos eletrônicos é um problema de classificação
binária com características que tornam sua solução não trivial:

\begin{itemize}[nosep]
  \item \textbf{Desbalanceamento extremo de classes.} A proporção de transações
    fraudulentas raramente ultrapassa 5\%, chegando a 0{,}17\% no dataset CC.
    Métricas clássicas como acurácia tornam-se desinformativas, exigindo o uso
    de ROC-AUC, PR-AUC, F1 e coeficiente de Matthews (MCC).

  \item \textbf{Deriva temporal (\textit{concept drift}).} Padrões de fraude
    evoluem continuamente; qualquer avaliação honesta deve respeitar a ordem
    cronológica das transações — treinando apenas com o passado e testando no
    futuro.

  \item \textbf{Heterogeneidade das features.} Os datasets modernos combinam
    variáveis numéricas contínuas, categóricas de alta cardinalidade, valores
    ausentes sistemáticos e features derivadas (agregações, codificações
    temporais).
\end{itemize}

Neste contexto, métodos baseados em \textit{gradient boosting} sobre árvores de
decisão (\algo{XGBoost}~\cite{chen2016xgboost}, \algo{LightGBM}~\cite{ke2017lightgbm},
\algo{CatBoost}~\cite{prokhorenkova2018catboost}) tornaram-se referência na
prática e em competições de \textit{machine learning}, pois lidam nativamente
com dados tabulares de alta dimensão, suportam valores ausentes, e possuem
mecanismos embutidos de ponderação por classe para desbalanceamento.

O objetivo deste trabalho é avaliar e comparar seis abordagens — incluindo um
método de detecção de anomalias (\algo{Isolation Forest}) e dois métodos de
ensemble de segundo nível (\algo{Stacking} e \algo{Super Ensemble}) — nos dois
datasets citados, utilizando métricas e protocolo experimental padronizados.

% =============================================================================
\section{Datasets}
\label{sec:datasets}
% =============================================================================

\subsection{Credit Card Fraud Detection (CC)}
\label{sec:dataset-cc}

O dataset \dataset{CC}~\cite{dal2015credit} foi disponibilizado pelo grupo de
pesquisa ULB Machine Learning Group e contém transações realizadas com cartões
de crédito europeus em setembro de 2013.

\begin{itemize}[nosep]
  \item \textbf{Dimensões:} \num{284807} transações $\times$ 31 colunas.
  \item \textbf{Taxa de fraude:} 0{,}17\% (492 fraudes em \num{284315} transações
    legítimas) — desbalanceamento severo de aproximadamente 1:578.
  \item \textbf{Features:} 28 componentes principais obtidas via PCA (V1--V28,
    anonimizadas para proteção de privacidade), o valor da transação (\texttt{Amount})
    e o tempo em segundos desde a primeira transação do dataset (\texttt{Time}).
\end{itemize}

A transformação PCA prévia elimina a multicolinearidade e anonimiza os dados,
mas também impede interpretação direta das features. O valor (\texttt{Amount})
é a única variável em escala original; sua distribuição é fortemente assimétrica
à direita, motivando transformações logarítmicas no pré-processamento do
\algo{Super Ensemble}.

\subsection{IEEE-CIS Fraud Detection (IEEE)}
\label{sec:dataset-ieee}

O dataset \dataset{IEEE}~\cite{ieee2019kaggle} foi disponibilizado pela Vesta
Corporation em parceria com a IEEE Computational Intelligence Society via
competição Kaggle (2019). É substancialmente mais complexo que o CC.

\begin{itemize}[nosep]
  \item \textbf{Dimensões:} \num{590540} transações $\times$ 434 colunas
    (resultado da junção das tabelas \texttt{train\_transaction} e
    \texttt{train\_identity} pela chave \texttt{TransactionID}).
  \item \textbf{Taxa de fraude:} 3{,}50\% (\num{20663} fraudes).
  \item \textbf{Features numéricas base:} valor da transação (\texttt{TransactionAmt}),
    14 features de contagem de comportamento do cartão (C1--C14), 5 features de
    distância temporal entre transações (D1--D5) e 20 features de identidade
    anônimas (V1--V20).
  \item \textbf{Features categóricas:} código do produto (\texttt{ProductCD}),
    bandeira e tipo do cartão (\texttt{card4}, \texttt{card6}) e domínios de
    e-mail do pagante e receptor (\texttt{P\_emaildomain}, \texttt{R\_emaildomain}).
  \item \textbf{Features adicionais no Super Ensemble:} 339 colunas V adicionais
    (filtradas por taxa de valores ausentes $<50\%$, resultando em 169 mantidas),
    variáveis de endereço, distância, indicadores de ausência para D1--D15,
    features de agregação por cartão e frequência de domínio de e-mail.
\end{itemize}

A Tabela~\ref{tab:datasets} resume as características principais de ambos os
datasets.

\begin{table}[H]
  \centering
  \caption{Características dos datasets utilizados nos experimentos.}
  \label{tab:datasets}
  \begin{tabular}{lrrrrr}
    \toprule
    \textbf{Dataset} & \textbf{Transações} & \textbf{Features} &
    \textbf{Fraudes} & \textbf{Taxa} & \textbf{Treino / Teste} \\
    \midrule
    CC   & \num{284807} & 29 (numéricas) & 492     & 0{,}17\% &
      \num{199364} / \num{42722} \\
    IEEE & \num{590540} & 40 (base)      & \num{20663} & 3{,}50\% &
      \num{413378} / \num{88581} \\
    \bottomrule
  \end{tabular}
\end{table}

% =============================================================================
\section{Metodologia}
\label{sec:metodologia}
% =============================================================================

\subsection{Protocolo de Divisão Treino/Teste}
\label{sec:split}

Para respeitar a causalidade temporal e evitar \textit{data leakage} de
informação futura, todos os modelos utilizam uma \textbf{divisão cronológica}
70\,/\,15\,/\,15 sobre a coluna de tempo de cada dataset (\texttt{Time} no CC
e \texttt{TransactionDT} no IEEE). O procedimento consiste em:

\begin{enumerate}[nosep]
  \item Ordenar todas as transações pelo timestamp.
  \item Atribuir os 70\% iniciais ao conjunto de \textbf{treino}.
  \item Atribuir os 15\% seguintes à \textbf{validação} (usada internamente
    por CatBoost e pelo early stopping dos demais modelos).
  \item Atribuir os 15\% finais ao conjunto de \textbf{teste} (avaliação final).
\end{enumerate}

Essa abordagem simula fielmente o cenário de implantação real, onde o modelo
aprende com transações passadas e é avaliado em transações futuras ainda não
vistas.

\subsection{Pré-processamento}
\label{sec:preprocessing}

O pré-processamento é diferenciado por família de método:

\subsubsection{Pré-processamento Base (XGBoost, LightGBM, CatBoost, Stacking)}

\begin{itemize}[nosep]
  \item \textbf{Imputação:} valores ausentes substituídos pela mediana calculada
    exclusivamente nos dados de treino (\texttt{SimpleImputer}).
  \item \textbf{Features numéricas:} padronização com \texttt{StandardScaler}
    (média 0, desvio 1) para o Isolation Forest; sem escalonamento para os
    modelos de árvore, que são invariantes a transformações monotônicas.
  \item \textbf{Features categóricas (IEEE):} codificação ordinal
    (\texttt{OrdinalEncoder}) com valor $-1$ para categorias desconhecidas no
    conjunto de teste; índices das colunas categóricas informados ao CatBoost
    para tratamento nativo.
  \item \textbf{Filtro de colunas V no IEEE:} das 339 colunas V disponíveis,
    apenas as 169 com taxa de valores ausentes inferior a 50\% (calculada no
    treino) são retidas, reduzindo ruído e custo computacional.
\end{itemize}

\subsubsection{Pré-processamento Enriquecido (Super Ensemble)}

O Super Ensemble recebe um espaço de features expandido, construído sem
\textit{data leakage} (todos os agregadores são ajustados exclusivamente no
conjunto de treino):

\paragraph{Dataset CC (40 features):}
\begin{itemize}[nosep]
  \item V1--V28 imputados pela mediana.
  \item \textbf{Normas por transação:} $\ell_2$ das sub-sequências V1--V14,
    V15--V28 e V1--V28 completo.
  \item \textbf{Estatísticas por transação:} média, desvio padrão, máximo,
    mínimo e contagem de valores negativos entre V1--V28.
  \item \textbf{Amount:} log\,($1 + \text{Amount}$) e escalonamento robusto
    (\texttt{RobustScaler}), que é insensível a outliers de fraude.
  \item \textbf{Tempo cíclico:} $\sin(2\pi h/24)$ e $\cos(2\pi h/24)$, onde
    $h$ é a hora do dia derivada de \texttt{Time}, capturando a periodicidade
    circadiana sem descontinuidade.
\end{itemize}

\paragraph{Dataset IEEE (247 features):}
Além de todas as features numéricas e categóricas expandidas:
\begin{itemize}[nosep]
  \item \textbf{Indicadores de ausência:} variáveis binárias para cada coluna
    D1--D15, pois a ausência em si pode ser informativa (distância temporal não
    registrada pode indicar primeira transação).
  \item \textbf{Features de agregação por cartão (\texttt{card1}):} contagem de
    transações, valor médio, desvio padrão e desvio normalizado do valor atual
    em relação ao histórico do cartão.
  \item \textbf{Features de agregação por endereço (\texttt{addr1}):} contagem
    e valor médio das transações.
  \item \textbf{Frequências de domínio de e-mail:} proporção de ocorrências de
    cada domínio (\texttt{P\_emaildomain}, \texttt{R\_emaildomain}) no treino.
  \item \textbf{Tempo cíclico e dia da semana:} codificação de \texttt{TransactionDT}
    em $\sin/\cos$ da hora e posição no ciclo semanal.
\end{itemize}

\subsection{Métricas de Avaliação}
\label{sec:metricas}

Dado o severo desbalanceamento de classes, o desempenho é reportado por um
conjunto amplo de métricas complementares:

\begin{description}[style=nextline]
  \item[\metric{ROC-AUC}]
    Área sob a curva ROC. Mede a capacidade discriminativa geral do modelo
    (probabilidade de ranquear uma fraude acima de uma transação legítima
    aleatória). Insensível ao limiar de decisão mas pode ser otimista em
    cenários de desbalanceamento extremo.

  \item[\metric{PR-AUC} (Average Precision)]
    Área sob a curva Precision-Recall. Foca no desempenho sobre a classe
    minoritária e é mais informativa que ROC-AUC quando o desbalanceamento é
    severo. Um classificador aleatório obteria PR-AUC $\approx$ taxa de fraude.

  \item[\metric{Gini}]
    $\text{Gini} = 2 \times \text{ROC-AUC} - 1$. Normalização da ROC-AUC no
    intervalo $[-1, 1]$, amplamente usada no setor financeiro.

  \item[\metric{KS} (Kolmogorov--Smirnov)]
    Máxima separação entre as distribuições acumuladas dos scores de fraude e
    de transações legítimas. Indica em que ponto do ranking o modelo melhor
    discrimina as classes.

  \item[\metric{MCC} (Matthews Correlation Coefficient)]
    Correlação entre classes reais e preditas. Considerada a métrica mais
    equilibrada para classificação binária desbalanceada, pois leva em conta
    todos os quadrantes da matriz de confusão.

  \item[\metric{F1}, \metric{Precision}, \metric{Recall}]
    Calculados no limiar de decisão $\tau^*$ que maximiza o F1 sobre o conjunto
    de teste — identificado via curva Precision-Recall discreta.

  \item[\metric{FPR@90R}]
    Taxa de falsos positivos (proporção de legítimas classificadas como fraude)
    no ponto da curva ROC em que o recall (sensibilidade) é de 90\%. Métrica
    operacional relevante: indica quantas transações legítimas precisam ser
    bloqueadas para capturar 90\% das fraudes.
\end{description}

% =============================================================================
\section{Métodos}
\label{sec:metodos}
% =============================================================================

\subsection{XGBoost}
\label{sec:xgboost}

\algo{XGBoost}~\cite{chen2016xgboost} (\textit{eXtreme Gradient Boosting}) é
um framework de \textit{gradient boosting} sobre árvores de decisão que combina
regularização $\ell_1/\ell_2$, subsampling de instâncias e features, e um
algoritmo de construção de árvores baseado em histogramas (\texttt{tree\_method=hist})
que reduz drasticamente o custo computacional em datasets grandes.

\paragraph{Configuração:}
\begin{itemize}[nosep]
  \item \texttt{n\_estimators=300}, com early stopping em 20 rodadas avaliando
    PR-AUC em um conjunto de validação interno (15\% do treino, estratificado).
  \item \texttt{scale\_pos\_weight} = $n_\text{neg}/n_\text{pos}$: ponderação
    que compensa o desbalanceamento ampliando o gradiente das amostras positivas.
  \item \texttt{eval\_metric=aucpr}: a otimização durante o early stopping
    prioriza a precisão sobre a classe de fraude, não a acurácia global.
\end{itemize}

\paragraph{Complexidade de treinamento:} $O(KDnF)$, onde $K$ é o número de
árvores, $D$ a profundidade máxima, $n$ o número de instâncias e $F$ o número
de features.

\subsection{LightGBM (DART)}
\label{sec:lightgbm}

\algo{LightGBM}~\cite{ke2017lightgbm} (\textit{Light Gradient Boosting Machine})
emprega crescimento de árvore \textit{leaf-wise} em vez de \textit{level-wise},
o que permite maior expressividade com o mesmo número de folhas. Neste
experimento foi utilizado o booster \textbf{DART} (\textit{Dropouts meet
Multiple Additive Regression Trees})~\cite{rashmi2015dart}, que aplica dropout
nas árvores existentes a cada iteração, com o objetivo de reduzir o
\textit{overfitting} e melhorar a generalização.

\paragraph{Configuração:}
\begin{itemize}[nosep]
  \item \texttt{boosting\_type=dart}, \texttt{drop\_rate=0.1}.
  \item \texttt{n\_estimators=1000} fixo (DART não suporta early stopping).
  \item \texttt{learning\_rate=0.05}, \texttt{num\_leaves=63},
    \texttt{subsample=0.8}, \texttt{colsample\_bytree=0.8}.
  \item \texttt{scale\_pos\_weight} igual ao XGBoost.
\end{itemize}

\paragraph{Nota sobre os resultados:} O LightGBM com DART apresentou colapso
de desempenho no dataset CC (ROC-AUC = 0{,}5037), provavelmente por
instabilidade do DART em datasets com desbalanceamento extremo (1:578) sem
early stopping. No dataset IEEE, com menor desbalanceamento (1:27), o método
operou normalmente (ROC-AUC = 0{,}8797).

\subsection{CatBoost}
\label{sec:catboost}

\algo{CatBoost}~\cite{prokhorenkova2018catboost} (\textit{Categorical Boosting})
introduz dois avanços principais: (1) ordered boosting, que evita o target
leakage durante o treinamento ao usar uma permutação aleatória das instâncias;
e (2) suporte nativo a features categóricas via combinações de target statistics
de segunda ordem, sem necessidade de one-hot encoding.

\paragraph{Configuração:}
\begin{itemize}[nosep]
  \item \texttt{iterations=1000}, \texttt{depth=8}, \texttt{learning\_rate=0.05}.
  \item \texttt{auto\_class\_weights=Balanced}: pesos automaticamente calculados
    como inverso da frequência de cada classe.
  \item \texttt{eval\_metric=AUC} com validação interna (15\% do treino).
  \item Índices das colunas categóricas passados explicitamente
    (\texttt{cat\_features}): 0 colunas no CC (todas PCA), 5 no IEEE.
\end{itemize}

O CatBoost foi o único modelo a exibir melhora monotônica no conjunto de
validação do IEEE ao longo das 1\,000 iterações completas (melhor iteração = 999),
indicando que o modelo ainda não havia saturado.

\subsection{Isolation Forest}
\label{sec:isoforest}

\algo{Isolation Forest}~\cite{liu2008isolation} é um método de detecção de
anomalias \textbf{não supervisionado} que parte da premissa de que instâncias
anômalas (fraudes) são raras e diferentes da maioria, sendo portanto mais
fáceis de ``isolar'' por particionamentos aleatórios do espaço de features.

\paragraph{Algoritmo:}
\begin{enumerate}[nosep]
  \item Construir $T$ árvores de isolamento (\textit{isolation trees}), cada
    uma selecionando aleatoriamente uma feature e um ponto de corte.
  \item Para cada instância, registrar o número médio de particionamentos
    necessários para isolá-la.
  \item O score de anomalia é inversamente proporcional ao comprimento médio do
    caminho: instâncias isoladas com poucos cortes recebem score elevado.
\end{enumerate}

\paragraph{Configuração:}
\begin{itemize}[nosep]
  \item \texttt{n\_estimators=100}, sem acesso aos rótulos de classe.
  \item Score = $-\texttt{decision\_function}(x)$, invertido para que valores
    maiores indiquem maior anomalia.
\end{itemize}

Por ser não supervisionado, o Isolation Forest não utiliza o rótulo de classe
durante o treinamento, o que o torna aplicável em cenários de \textit{zero-shot}
e sem dados rotulados. Em contrapartida, sua PR-AUC é muito inferior (CC:
0{,}0434; IEEE: 0{,}0678), pois o modelo não distingue especificamente o padrão
de fraude do padrão geral de anomalia.

\subsection{Stacking}
\label{sec:stacking}

O \algo{Stacking}~\cite{wolpert1992stacked} é uma técnica de ensemble de
segundo nível em que as predições de múltiplos modelos base são usadas como
features para um meta-aprendiz. Para evitar \textit{data leakage}, as predições
para o treinamento do meta-aprendiz são geradas por validação cruzada
\textit{out-of-fold} (OOF).

\paragraph{Arquitetura:}
\begin{itemize}[nosep]
  \item \textbf{Modelos base:} XGBoost, LightGBM (GBDT), CatBoost, Regressão
    Logística (com imputação e regularização $\ell_2$, $C=0{,}1$).
  \item \textbf{Validação cruzada:} 5-fold estratificada.
  \item \textbf{Meta-aprendiz:} Regressão Logística ($C=1{,}0$) treinada nas
    predições OOF concatenadas ($n_\text{treino} \times 4$ features).
  \item \textbf{Inferência:} média das predições dos 5 modelos de cada fold,
    passada ao meta-aprendiz.
\end{itemize}

Os coeficientes do meta-aprendiz no CC foram: XGBoost = 5{,}062;
CatBoost = 2{,}749; LR = 3{,}245; LightGBM = 1{,}140 — refletindo a dominância
do XGBoost e CatBoost. No IEEE, LightGBM e LR receberam coeficientes negativos
ou próximos de zero, indicando que o meta-aprendiz os descartou efetivamente.

\subsection{Super Ensemble}
\label{sec:superensemble}

O \algo{Super Ensemble} estende o Stacking com três modificações: (1) features
enriquecidas com engenharia de atributos especializada; (2) modelos base com
hiperparâmetros mais agressivos; e (3) meta-aprendiz XGBoost otimizado por
PR-AUC, mais adequado ao desbalanceamento.

\paragraph{Arquitetura:}
\begin{itemize}[nosep]
  \item \textbf{Modelos base:} XGBoost (\texttt{n\_estimators=1000},
    \texttt{max\_depth=6}, regularização $\gamma/\alpha/\lambda$),
    LightGBM (\texttt{n\_estimators=2000}, \texttt{num\_leaves=63},
    \texttt{learning\_rate=0.03}),
    CatBoost (\texttt{iterations=2000}, \texttt{depth=8},
    \texttt{learning\_rate=0.03}).
  \item \textbf{Paralelismo:} os 3 modelos base são treinados simultaneamente
    via \texttt{ThreadPoolExecutor}, com os cores de CPU divididos igualmente.
  \item \textbf{Validação cruzada:} 5-fold estratificada (OOF).
  \item \textbf{Meta-features:} 7 features por instância — as 3 probabilidades
    brutas, 3 produtos par a par ($p_\text{xgb} \cdot p_\text{lgb}$,
    $p_\text{xgb} \cdot p_\text{cat}$, $p_\text{lgb} \cdot p_\text{cat}$) e
    o espalhamento $\max(\mathbf{p}) - \min(\mathbf{p})$.
  \item \textbf{Meta-aprendiz:} XGBoost com \texttt{eval\_metric=aucpr} e
    early stopping (30 rodadas) em 20\% do OOF.
  \item \textbf{Modelos finais:} re-treinados em 100\% do treino (com 10\% de
    holdout estratificado apenas para early stopping).
\end{itemize}

% =============================================================================
\section{Resultados}
\label{sec:resultados}
% =============================================================================

\subsection{Tabela de Resultados Consolidados}
\label{sec:tabela}

A Tabela~\ref{tab:resultados} apresenta os resultados completos de todos os
algoritmos nos dois datasets. Os melhores valores de cada métrica por dataset
estão destacados em \textbf{negrito}.

\begin{table}[H]
  \centering
  \caption{Resultados comparativos — todos os algoritmos nos datasets CC e IEEE.
  O limiar de decisão é selecionado por maximização do F1 no conjunto de teste.}
  \label{tab:resultados}
  \small
  \setlength{\tabcolsep}{4pt}
  \begin{tabular}{llrrrrrrrr}
    \toprule
    \textbf{Algoritmo} & \textbf{DS} &
    \textbf{ROC-AUC} & \textbf{PR-AUC} & \textbf{Gini} & \textbf{KS} &
    \textbf{MCC} & \textbf{F1} & \textbf{FPR@90R} & \textbf{t (s)} \\
    \midrule
    XGBoost          & CC & 0,9587 & 0,7402 & 0,9174 & 0,8254 & 0,8097 & 0,8000 & 0,0942 &    1,8 \\
    Isolation Forest & CC & 0,9459 & 0,0434 & 0,8918 & 0,7736 & 0,2037 & 0,1488 & 0,2077 &    0,5 \\
    LightGBM         & CC & 0,5037 & 0,0012 & 0,0073 & 0,0073 & 0,0024 & 0,0036 & 0,8993 &   32,8 \\
    CatBoost         & CC & \textbf{0,9713} & \textbf{0,7559} & \textbf{0,9425} & 0,8287 & 0,8097 & 0,8000 & \textbf{0,0808} &   16,4 \\
    Stacking         & CC & 0,9643 & 0,7513 & 0,9286 & \textbf{0,8319} & \textbf{0,8246} & \textbf{0,8211} & 0,1011 &   89,3 \\
    Super Ensemble   & CC & 0,9290 & 0,6994 & 0,8580 & 0,8155 & 0,8214 & 0,8132 & 0,1325 &   61,8 \\
    \midrule
    XGBoost          & IEEE & 0,8443 & 0,4462 & 0,6886 & 0,5400 & 0,4735 & 0,4729 & 0,5415 &    6,8 \\
    Isolation Forest & IEEE & 0,6079 & 0,0678 & 0,2158 & 0,1877 & 0,0775 & 0,1156 & 0,8453 &    0,6 \\
    LightGBM         & IEEE & 0,8797 & 0,4719 & 0,7594 & 0,6055 & 0,4769 & 0,4752 & 0,4093 &  210,1 \\
    CatBoost         & IEEE & 0,8687 & 0,4731 & 0,7375 & 0,5902 & 0,4814 & 0,4876 & 0,4663 &  186,3 \\
    Stacking         & IEEE & 0,8540 & 0,4482 & 0,7080 & 0,5644 & 0,4683 & 0,4691 & 0,4967 &  266,4 \\
    Super Ensemble   & IEEE & \textbf{0,8946} & \textbf{0,5095} & \textbf{0,7891} & \textbf{0,6241} & \textbf{0,5194} & \textbf{0,5255} & \textbf{0,3556} & 1729,9 \\
    \bottomrule
  \end{tabular}
\end{table}

\subsection{Análise por Dataset}
\label{sec:analise}

\subsubsection{Dataset CC}

O dataset CC apresenta características que favorecem fortemente os métodos
supervisionados de \textit{gradient boosting}:

\begin{itemize}[nosep]
  \item \textbf{CatBoost} obteve o melhor ROC-AUC (0{,}9713) e PR-AUC
    (0{,}7559), além da menor FPR@90R (0{,}0808), indicando que, para capturar
    90\% das fraudes, apenas 8\% das transações legítimas são erroneamente
    bloqueadas. Isso se deve à eficácia do \textit{ordered boosting} e à
    ponderação automática de classes.

  \item \textbf{Stacking} obteve os melhores F1 (0{,}8211), MCC (0{,}8246) e
    KS (0{,}8319), superando o CatBoost nestas métricas ao combinar
    complementariedade dos modelos base via meta-aprendiz.

  \item \textbf{LightGBM (DART)} colapsou completamente no CC (ROC-AUC =
    0{,}5037, equivalente a classificação aleatória). O uso do DART sem early
    stopping em dados com desbalanceamento 1:578 resultou em instabilidade
    numérica — o modelo convergiu para uma solução degenerada que essencialmente
    prediz classe negativa para todas as instâncias.

  \item \textbf{Isolation Forest} demonstra boa capacidade discriminativa em
    termos de ROC-AUC (0{,}9459), mas PR-AUC ínfima (0{,}0434), evidenciando
    que, embora ranqueie corretamente as fraudes em posição elevada, a densidade
    de fraudes entre os top scores é muito baixa — o modelo confunde outras
    anomalias (valores extremos de Amount, por exemplo) com fraudes.

  \item \textbf{Super Ensemble} apresentou resultado inferior ao Stacking e ao
    CatBoost isolado no CC, sugerindo que o enriquecimento de features adiciona
    mais ruído do que sinal neste dataset, cujas features PCA já são altamente
    informativas.
\end{itemize}

\subsubsection{Dataset IEEE}

O dataset IEEE, com maior complexidade e menor desbalanceamento relativo,
apresenta ordenamento diferente:

\begin{itemize}[nosep]
  \item \textbf{Super Ensemble} domina em todas as métricas principais:
    ROC-AUC = 0{,}8946, PR-AUC = 0{,}5095, F1 = 0{,}5255 e a menor FPR@90R
    (0{,}3556). O enriquecimento de features (features de agregação por cartão,
    indicadores de ausência, tempo cíclico) foi determinante, pois o IEEE possui
    variáveis com semântica clara (diferentemente das PCA do CC).

  \item \textbf{LightGBM} recupera desempenho no IEEE (ROC-AUC = 0{,}8797),
    confirmando que o problema observado no CC foi específico ao desbalanceamento
    extremo com DART.

  \item \textbf{Stacking} não supera os modelos base isolados no IEEE, com
    ROC-AUC = 0{,}8540 inferior ao XGBoost (0{,}8443), LightGBM (0{,}8797) e
    CatBoost (0{,}8687). O meta-aprendiz atribuiu coeficientes negativos ao
    LightGBM e à LR, descartando-os efetivamente.

  \item \textbf{XGBoost isolado} é o algoritmo mais rápido entre os supervisionados
    (6{,}8s no IEEE), obtendo ROC-AUC = 0{,}8443 — competitivo considerando a
    simplicidade e velocidade.
\end{itemize}

\subsection{Análise das Curvas ROC e Precision-Recall}
\label{sec:curvas}

As Figuras~\ref{fig:roc} e~\ref{fig:pr} mostram as curvas ROC e
Precision-Recall para todos os algoritmos nos dois datasets.

\begin{figure}[H]
  \centering
  \includegraphics[width=\textwidth]{plots/results/01_roc_curves.png}
  \caption{Curvas ROC por dataset. A diagonal tracejada representa um
  classificador aleatório (AUC = 0{,}5).}
  \label{fig:roc}
\end{figure}

\begin{figure}[H]
  \centering
  \includegraphics[width=\textwidth]{plots/results/02_pr_curves.png}
  \caption{Curvas Precision-Recall por dataset. O nível de precisão de um
  classificador aleatório corresponde à taxa de fraude do dataset (0{,}17\%
  no CC, 3{,}50\% no IEEE).}
  \label{fig:pr}
\end{figure}

A divergência entre ROC-AUC e PR-AUC do Isolation Forest é visualmente
notável: a curva ROC aparenta desempenho razoável, enquanto a curva PR revela
que o modelo falha em identificar fraudes com alta precisão — demonstrando a
importância de avaliar ambas as curvas em problemas desbalanceados.

\subsection{Distribuição dos Scores e Separabilidade}
\label{sec:scores}

\begin{figure}[H]
  \centering
  \includegraphics[width=\textwidth]{plots/results/05_score_distributions.png}
  \caption{Distribuições dos scores de anomalia por classe (fraude vs.\ legítima),
  recortadas nos percentis 1--99 para legibilidade.}
  \label{fig:scores}
\end{figure}

\begin{figure}[H]
  \centering
  \includegraphics[width=\textwidth]{plots/results/06_ks_curves.png}
  \caption{Curvas KS — distribuições cumulativas dos scores por classe. A barra
  vertical indica o ponto de máxima separação (estatística KS).}
  \label{fig:ks}
\end{figure}

A estatística KS reflete diretamente a separabilidade entre as distribuições
de scores das duas classes. No CC, o CatBoost obteve KS = 0{,}8287, enquanto o
Stacking atingiu KS = 0{,}8319 — indicando que ambos os modelos concentram as
fraudes no extremo superior do ranking com eficácia similar. No IEEE, o Super
Ensemble apresentou KS = 0{,}6241, o maior dentre todos os métodos avaliados.

\subsection{Análise de Threshold e Matrizes de Confusão}
\label{sec:threshold}

\begin{figure}[H]
  \centering
  \includegraphics[width=\textwidth]{plots/results/08_threshold_analysis.png}
  \caption{Precision, Recall e F1 em função do threshold de decisão. A linha
  tracejada marca o threshold ótimo $\tau^*$ (maximização do F1).}
  \label{fig:threshold}
\end{figure}

\begin{figure}[H]
  \centering
  \includegraphics[width=\textwidth]{plots/results/09_confusion_matrices.png}
  \caption{Matrizes de confusão normalizadas por linha (taxa real) no threshold
  ótimo $\tau^*$.}
  \label{fig:confusion}
\end{figure}

As matrizes de confusão evidenciam o trade-off operacional: no CC, XGBoost e
CatBoost obtêm Precision = 0{,}9474 ao custo de Recall = 0{,}6923 — ou seja,
identificam apenas 69\% das fraudes, mas com altíssima precisão nas alertas.
O Stacking equilibra melhor este trade-off (Recall = 0{,}7500, Precision =
0{,}9070). A escolha do threshold ideal depende do custo relativo de falsos
negativos (fraude não detectada) versus falsos positivos (transação legítima
bloqueada) — parâmetro definido pela política de negócio de cada organização.

\subsection{Curvas de Ganho Cumulativo}
\label{sec:ganho}

\begin{figure}[H]
  \centering
  \includegraphics[width=\textwidth]{plots/results/07_gain_curves.png}
  \caption{Curvas de ganho cumulativo: fração de fraudes capturadas em função
  da fração de transações inspecionadas (ranking por score decrescente).}
  \label{fig:gain}
\end{figure}

As curvas de ganho cumulativo ilustram o valor prático do modelo em um cenário
de inspeção manual: inspecionando apenas os 10\% de transações com maior score,
o CatBoost captura aproximadamente 90\% das fraudes no CC. No IEEE, o Super
Ensemble captura cerca de 70\% das fraudes inspecionando os 10\% de maior score.

\subsection{Custo Computacional}
\label{sec:tempo}

\begin{figure}[H]
  \centering
  \includegraphics[width=\textwidth]{plots/results/04_training_times.png}
  \caption{Tempo de treinamento por algoritmo e dataset (escala linear).}
  \label{fig:tempo}
\end{figure}

A Tabela~\ref{tab:resultados} e a Figura~\ref{fig:tempo} revelam variação de
três ordens de magnitude no tempo de treinamento:

\begin{itemize}[nosep]
  \item \textbf{Isolation Forest} é o mais rápido ($<1$s em ambos os datasets),
    por ser não supervisionado e não requerer otimização de gradiente.
  \item \textbf{XGBoost} oferece o melhor equilíbrio custo-benefício:
    1{,}8s/6{,}8s com desempenho competitivo.
  \item \textbf{Super Ensemble} consumiu 1\,729{,}9s ($\approx$29 min) no IEEE,
    reflexo de 5-fold OOF com 3 modelos agressivos ($\leq$2\,000 iterações cada)
    mais re-treinamento final — justificado pelo ganho de desempenho.
\end{itemize}

% =============================================================================
\section{Discussão}
\label{sec:discussao}
% =============================================================================

\subsection{Impacto do Desbalanceamento}

O desbalanceamento é o fator mais determinante das diferenças de comportamento
observadas. A ponderação por classe (\texttt{scale\_pos\_weight} no XGBoost/LightGBM,
\texttt{auto\_class\_weights=Balanced} no CatBoost) é essencial para que o
modelo não colapse em predizer sempre a classe majoritária. A PR-AUC é
consistentemente mais baixa que a ROC-AUC em todos os modelos — e a razão entre
elas diminui conforme o desbalanceamento aumenta, como previsto teoricamente.

\subsection{DART e Instabilidade no CC}

O colapso do LightGBM com DART no CC é um resultado relevante. O DART utiliza
dropout sobre as árvores, o que pode causar alta variância nas estimativas de
gradiente em datasets com pouquíssimos positivos. Sem early stopping (DART não
o suporta), o modelo treinou por 1\,000 iterações em direção a um mínimo local
degenerado. A solução natural seria substituir o DART pelo booster padrão GBDT
com early stopping, como feito no Stacking.

\subsection{Generalização entre Datasets}

O Super Ensemble generaliza melhor para o IEEE, onde o enriquecimento de
features tem impacto significativo: as features de agregação por cartão
(contagem, desvio do valor em relação ao histórico) são proxies para o
comportamento habitual do portador — informação que não existe no CC, onde as
features PCA já capturam toda a estrutura relevante.

\subsection{Métricas Operacionais}

A FPR@90R é a métrica mais relevante para implantação em produção. No CC, o
CatBoost (FPR@90R = 0{,}0808) é superior ao Stacking (0{,}1011) e ao Super
Ensemble (0{,}1325), indicando que, para capturar 90\% das fraudes, o CatBoost
bloqueia apenas 8\% das transações legítimas — potencialmente $\sim$3\,400
transações bloqueadas por dia em um sistema de médio porte, um volume operacional
aceitável. No IEEE, o Super Ensemble reduz este valor para 35{,}6\%.

% =============================================================================
\section{Conclusão}
\label{sec:conclusao}
% =============================================================================

Este trabalho avaliou seis métodos baseados em árvores de decisão para detecção
de fraudes em dois benchmarks de referência. Os principais achados são:

\begin{enumerate}
  \item \textbf{CatBoost e XGBoost} são as melhores escolhas para implantação
    quando velocidade e interpretabilidade são prioritárias, oferecendo
    excelente desempenho com configuração direta.

  \item \textbf{Stacking} melhora consistentemente o F1 e o MCC ao combinar
    modelos complementares, mas não garante ganho em ROC-AUC — seu ponto forte
    é o equilíbrio entre Precision e Recall.

  \item \textbf{Super Ensemble} é a melhor opção quando o tempo de treinamento
    não é restrição e as features brutas têm semântica explorada por engenharia
    de atributos, como no IEEE.

  \item \textbf{Isolation Forest} é útil como \textit{baseline} não
    supervisionado ou para detecção de novas modalidades de fraude não
    rotuladas, mas não compete com métodos supervisionados em PR-AUC.

  \item \textbf{LightGBM com DART} mostrou-se frágil sob desbalanceamento
    extremo sem early stopping; o booster GBDT seria mais adequado neste cenário.

  \item A \textbf{divisão cronológica} é imprescindível para avaliação fidedigna:
    uma divisão aleatória inflacionaria artificialmente as métricas ao permitir
    que o modelo aprenda padrões de transações futuras.
\end{enumerate}

Como trabalhos futuros, sugere-se: (i) avaliar GBDT (sem DART) para o LightGBM
no CC; (ii) explorar calibração de probabilidade (Platt scaling, isotonic
regression) para melhorar a estimação das probabilidades de fraude; (iii)
investigar o impacto de hiperparâmetros específicos de PR-AUC nos modelos base
do Stacking.

% =============================================================================
%  Referências
% =============================================================================
\begin{thebibliography}{99}

\bibitem{chen2016xgboost}
T.~Chen e C.~Guestrin,
``XGBoost: A Scalable Tree Boosting System,''
\textit{Proc.\ ACM SIGKDD}, pp.~785--794, 2016.

\bibitem{ke2017lightgbm}
G.~Ke, Q.~Meng, T.~Finley, T.~Wang, W.~Chen, W.~Ma, Q.~Ye, e T.-Y.~Liu,
``LightGBM: A Highly Efficient Gradient Boosting Decision Tree,''
\textit{Advances in Neural Information Processing Systems (NeurIPS)}, 2017.

\bibitem{prokhorenkova2018catboost}
L.~Prokhorenkova, G.~Gusev, A.~Vorobev, A.~V.~Dorogush, e A.~Gulin,
``CatBoost: Unbiased Boosting with Categorical Features,''
\textit{Advances in Neural Information Processing Systems (NeurIPS)}, 2018.

\bibitem{liu2008isolation}
F.~T.~Liu, K.~M.~Ting, e Z.-H.~Zhou,
``Isolation Forest,''
\textit{Proc.\ IEEE ICDM}, pp.~413--422, 2008.

\bibitem{wolpert1992stacked}
D.~H.~Wolpert,
``Stacked Generalization,''
\textit{Neural Networks}, vol.~5, no.~2, pp.~241--259, 1992.

\bibitem{rashmi2015dart}
K.~V.~Rashmi e R.~Gilad-Bachrach,
``DART: Dropouts Meet Multiple Additive Regression Trees,''
\textit{Proc.\ AISTATS}, 2015.

\bibitem{dal2015credit}
A.~Dal Pozzolo, O.~Caelen, R.~A.~Johnson, e G.~Bontempi,
``Calibrating Probability with Undersampling for Unbalanced Classification,''
\textit{Proc.\ IEEE SSCI}, 2015.

\bibitem{ieee2019kaggle}
Vesta Corporation e IEEE-CIS,
``IEEE-CIS Fraud Detection,''
\textit{Kaggle Competition}, 2019.
\url{https://www.kaggle.com/c/ieee-fraud-detection}

\end{thebibliography}

\end{document}

